\documentclass{article}
\usepackage[utf8]{inputenc}
\usepackage{graphicx}
\usepackage{amsmath}
%\usepackage[fceqn]{amsmath,amssymb}
\usepackage[bottom=2.0cm,top=2.0cm,left=2.5cm,right=2.5cm]{geometry}
\usepackage[portuges]{babel}
\usepackage{indentfirst}
\usepackage{wrapfig}
\usepackage{ mathrsfs }
\usepackage{subcaption}
\usepackage{ textcomp }
\usepackage[export]{adjustbox}
\usepackage{setspace}
\usepackage{stackengine}
\usepackage{hyperref}  %%%%
\setlength\parindent{0pt}
\hypersetup{colorlinks,citecolor=black,filecolor=black,linkcolor=black,urlcolor=black} %%%%




%\title{Atividades 4 e 5}
%\author{Jéferson André Sales Fortunato}

\begin{document}



%\begin{titlepage} put this to make a tittle page
\begin{center}
    \includegraphics[scale=0.13, valign=c]{images/unifesp.png}
\end{center}
\begin{minipage}[c]{1\textwidth}
\centering
\large{ {\bfseries UNIVERSIDADE FEDERAL DE SÃO PAULO - UNIFESP} \\ {\bfseries DEPARTAMENTO DE FÍSICA} \\ {\bfseries DISCIPLINA}: Disciplina associada à atividade \\ {\bfseries PROFESSOR:} Nome do/da professor/a \\ {\bfseries ESTUDANTE:} Seu nome }
\end{minipage}

%\end{titlepage}

%\maketitle

\begin{large}

\section*{Atividade X}

\textbf{Questão X:} Enunciado da questão X.
\\

\noindent\textbf{Solução:}
\\

\begin{equation}
    A+B=C+D.
\end{equation}
\\

\textbf{Questão X:} Enunciado da questão X.
\\

\noindent\textbf{Solução:}
\\

\begin{equation}
    A+B=C+D.
\end{equation}
\\

\textbf{Questão X:} Enunciado da questão X.
\\

\noindent\textbf{Solução:}
\\

\begin{equation}
    A+B=C+D.
\end{equation}
\\

\textbf{Questão X:} Enunciado da questão X.
\\

\noindent\textbf{Solução:}
\\

\begin{equation}
    A+B=C+D.
\end{equation}
\\

\end{large}
\end{document}